\section{Discussion}\label{sec:discussion}

\statspai{} is a unification and validation project, not a new-estimator
paper. The claim is that a single \proglang{Python} package can expose
the applied causal workflow through shared result surfaces, explicit
validation status, reproducible parity artifacts, and machine-readable
agent schemas. It is not a claim that \statspai{} dominates specialized
reference packages in their home domains, nor that the MCP interface
replaces empirical judgement.

The main limitations are as follows. First, \RegistryAutoUnbacked{} stable auto-registered symbols remain API-stable but not parity-backed
(\RegistryAutoUnbackedClasslike{} class-like, \RegistryAutoUnbackedFunctionlike{} function-like), and the other \RegistryOtherApiEvidence{} symbols of the
\RegistryApiStable-symbol \code{api\_stable} tier carry API, unit-contract,
convention, or supplemental evidence only. The validated claim is
limited to certified/validated entries and to the named rows and modules in the validation-suite evidence; the validation-evidence
audit lists nine scoped limitation rows, each tied to a named function, evidence grade, and unsupported variant.
Second, only one Track~A row remains a methodological/T4 disclosure:
classical SCM on the calibrated Basque replica, where reference
implementations choose measurably different local optima; on the
original \code{Synth::basque} bytes the same specification brings all
three implementations within \(4\times10^{-4}\)
(Section~\ref{sec:scm-identification}), so the disclosure is about
data-dependent non-uniqueness, not a standing numerical defect. The
causal-forest row is T3: the two forest engines are equivalent to
within a tenth of a sampling standard error and equally noisy across
seeds, which is an equivalence claim, not parity. The
\ParityThirdPartyModuleCount{} Track~A rows served by third-party
\proglang{Python} libraries validate \statspai{}'s wrappers of those
libraries, not native algorithms. The configuration map behind
\code{sp.validation\_scope} covers the twelve suite estimators; for
the rest of the registry the evidence notes name the artifact but not
its configuration. Third,
the \pkg{fect} and \pkg{interflex} cross-validation selectors for the
factor number, penalty, and kernel bandwidth are ported in
1.29.0 but can only be graded T3: their folds are random,
so the three-ecosystem parity of Section~\ref{sec:ex-fect} is
established at user-specified tuning values, and the selectors are
checked by given-fold score equality plus across-seed agreement of the
selection. Fourth, Stata reproduction requires a
separate licence, so Stata rows are migration bridges rather than the
licence-free core of the archive; \NoStataModuleCount{} of the \RParityModuleCount{} modules carry no Stata
reference, each with a measured reason. Fifth, Bayesian, causal-discovery,
and behavioural-agent evaluations require different loss functions and
are outside this software-validation article. Sixth, within DiD, \statspai{} implements the apparatus of the
\citet{baker2026difference} practitioner's guide but not the efficient
combination of comparison groups and baseline periods
\citep{chen2025efficient} or selection-based pre-trend diagnostics
\citep{ghanem2026selection}.

The unification itself has costs worth stating plainly, because they
are the price of the design rather than defects of it. A unified API
necessarily lags its upstream references: when \pkg{did} or
\pkg{fixest} ships a new option, \statspai{} users wait for the ported,
parity-tested version, and the validation regime that makes the port
trustworthy is exactly what makes it slower to ship. The registry's
own census quantifies a second cost: \RegistryAutoUnbacked{} of the
registered symbols carry API-contract evidence only, so a user who
strays from the validated core onto a long-tail helper is relying on
software-engineering guarantees, not statistical validation -- which is
why the tier is printed at call time rather than in an appendix.
Third, the evidence regime itself has a maintenance bill: every
correctness fix in Table~\ref{tab:correctness-history} exists because a
parity or coverage artifact flagged it, and keeping \RParityModuleCount{}
modules pinned against two external ecosystems is ongoing work that a
single-method package does not carry.

A fourth cost is a limit of the regime rather than of the package. A
parity ledger validates the paths it exercises and is silent about the
rest: the 1.25.0 entry of Table~\ref{tab:correctness-history}, a
continuous-treatment forest path no Track~A or Track~B row reaches,
surfaced only in a worked example. That is why the registry records
\emph{which} artifact backs a label, why \code{sp.validation\_scope}
reports the configuration an artifact exercised, and why the worked
examples are executed end-to-end.

The roadmap is to resolve or further bound the remaining
classical-SCM reference-disagreement disclosure, extend
\(B=1{,}000\) coverage beyond the headline rows, close the two measured
performance gaps of Section~\ref{sec:performance}, and move the agent
interface from mechanical conformance to scored behavioural
benchmarking.

A final limitation is institutional: a large surface is maintained by a
small team with one principal author. The risk is bounded by
properties verifiable from the archive: the package is MIT-licensed with
no relicensing clause, so the validated core can be forked; every
output-changing fix is pinned by a regression test; and the parity
harness is scripted and environment-pinned (\code{renv.lock}, recorded
Stata versions, per-row JSON provenance), so any third party with the
reference runtimes can re-execute it.

\section*{Author contributions}

B.\,W.\ designed and implemented \statspai{}, the validation harnesses
and registry API, and drafted the manuscript. S.\,R.\ supervised the
applied-econometrics programme, advised on scope and validation priorities,
and reviewed the manuscript. Both authors approved the final version.

\section*{Acknowledgements}

The authors thank Stanford REAP for institutional support and feedback,
and the maintainers of the \proglang{Python}, \proglang{R}, and
\proglang{Stata} packages used as parity references. Any remaining
errors are the authors' own.

\section*{Conflict of interest}

\statspai{} is MIT-licensed open-source software. The authors are
affiliated with StatsPAI Inc., which maintains the package; CoPaper.AI
is a commercial downstream product that may call \statspai{} but is not
required to use it. B.\,W.\ is the creator and principal maintainer. No
external funder had a role in the validation studies, analysis, or
manuscript preparation.

\section*{AI usage disclosure}

Portions of code documentation and manuscript text were prepared with
large-language-model assistance. The authors selected the claims,
reviewed the text, and remain responsible for all statistical decisions.
Numerical claims are tied to tests, replication scripts, frozen JSON
artifacts, and cited references; citations are checked by the verifier
in Section~\ref{sec:cite}.
